\chapter{Excess-Flow Valves in Hydrogen refueling}
\label{ch:efv_operation}

\section{Hydrogen Refueling Infrastructure and Safety Functions}
\label{sec:Refueling_context}

The design of an excess-flow valve (EFV) needs to be considered in the context of the system in which it operates. A gaseous hydrogen Refueling station typically brings together a hydrogen supply, compression and buffer-storage equipment, temperature-control devices, and a dispensing unit connected to the vehicle through a hose and nozzle, as the ISO~19880-1 indicates \cite{iso19880_1}. The station must deliver the intended flow during a normal fill while maintaining the ability to limit or isolate releases after a malfunction.


\begin{figure}[htbp]
\centering
\resizebox{\textwidth}{!}{% Conceptual equipment sequence, not a mandatory ISO layout.
\begin{tikzpicture}[
  >=Latex, line width=0.8pt,
  item/.style={draw,rounded corners=2pt,fill=black!3,align=center,
               minimum height=1.2cm,font=\small,inner sep=4pt},
  flow/.style={->,line width=1pt},
  note/.style={font=\footnotesize,align=center}
]
\node[item,minimum width=1.9cm] (s) at (0,0) {Hydrogen\\supply};
\node[item,minimum width=2.7cm] (b) at (3.1,0) {Compression and\\buffer storage};
\node[item,minimum width=2.0cm] (c) at (6.2,0) {Pre-cooler};
\node[item,minimum width=1.65cm,thick,fill=black!10] (v) at (8.9,0) {\textbf{EFV}};
\node[item,minimum width=2.05cm] (d) at (11.6,0) {Dispenser};
\node[item,minimum width=2.0cm] (t) at (14.5,0) {Vehicle\\tank};
\draw[flow] (s.east)--(b.west);
\draw[flow] (b.east)--(c.west);
\draw[flow] (c.east)--(v.west);
\draw[flow] (v.east)--(d.west);
\draw[flow] (d.east)--(t.west);
\node[note] at (8.9,1.15) {Selected EFV position\\after pre-cooling};
\draw[->] (8.9,0.77)--(v.north);
\node[note] at (7.45,-1.20) {Pre-cooled, high-pressure hydrogen};
\draw[densely dashed] (5.1,-0.83)--(10.1,-0.83);
\node[note] at (8.9,-1.84) {Reference state for the proposed operating study:\\
$p_{\mathrm{ref}}=430\,\mathrm{bar}$ (absolute),\quad $T_{\mathrm{ref}}=-40\,^{\circ}\mathrm{C}$};
\node[note] at (14.5,1.08) {H35/H70\\storage class};
\end{tikzpicture}
}
\caption{Hydrogen refueling system flow path scheme with highlighted location of the EFV}
\label{fig:station_efv_location}
\end{figure}

For normal duty vehicles refueling, SAE~J2601 specifies fueling protocols for the H35 and H70 nominal pressure classes, corresponding to $35$ and $70\,\mathrm{MPa}$, and identifies fuel delivery temperature categories T20, T30 and T40. The T40 designation refers to a nominal fuel-delivery temperature of $-40\,^{\circ}\mathrm{C}$. The fueling protocol constrains also quantities such as mass-flow rate, pressure-rise rate, delivery temperature and end pressure, these values will be used as references for the successive CFD simulation \cite{sae_j2601}. Gas pressure and temperature change during a filling event, this requires to specify which system section is studied.

An EFV is just one element of a wider protective arrangement, alongside check valves, flow-control valves, breakaway devices, pressure safety valves and actuated shut-off valves \cite{iso19880_3}. 
In particular the purpose of an EFV is limiting abnormal \emph{forward} flow; while check valves prevent reverse flow and emergency shut-off, instead an hose breakaway addresses the specific mechanical separation of the dispenser connection. These functions are complementary and should not be represented as interchangeable \cite{iso19880_1,iso19880_3,h2tools_valves}.

From a component design perspective, the relevant questions are not limited to whether a valve can shut off flow. The designer must also consider the pressure loss in normal service, the repeatability of its response, the influence of gas properties, and whether the required hydrodynamic force can be generated at the intended excess-flow condition. These competing requirements motivate the geometrical and fluid-dynamic comparisons developed in the following chapters.

\section{EFV Installation Point and Reference Operating State}
\label{sec:efv_reference_state}

For the industrial application addressed in this thesis, the excess-flow valve is located \emph{downstream of the pre-cooling heat exchanger, in the high-pressure delivery line close to the dispenser}, as indicated in Figure~\ref{fig:station_efv_location}, this is not a normative required location. The standard only establishes requirements and test methods for high-pressure hydrogen-station valves, including excess-flow valves, and distinguishes several valve functions \cite{iso19880_3}.

Pre-cooling is a relevant aspect for fast vehicle refueling because compression of the gas in the on-board storage system raises its temperature during filling. The normative SAE~J2601 addresses H35 ($35\,\mathrm{MPa}$ nominal working pressure) and H70 ($70\,\mathrm{MPa}$) vehicle storage and defines the T40 delivery-temperature category. Pre-cooling stations are usually arranged so that a heat exchanger conditions the hydrogen before it passes to the dispenser \cite{sae_j2601,elgowainy2017,iso19880_1}. Therefore, ambient Hydrogen properties are not representative of the selected valve's intended cold delivery-line condition.

A representative reference state for the pursued investigation is therefore selected as:
\begin{equation}
  p_{\mathrm{ref}}=437\,\mathrm{bar}=43.7\,\mathrm{MPa},
  \qquad
  T_{\mathrm{ref}}\simeq-40\,^{\circ}\mathrm{C}=233\,\mathrm{K}.
  \label{eq:efv_reference_state}
\end{equation}

The selected pair represents a specified design operating point, \emph{not} a pressure--temperature requirement universally imposed by SAE~J2601 or a claim that pressure remains constant throughout an entire fill. In particular, $437\,\mathrm{bar}$ should not be mistaken for the complete operating envelope of an H70 dispenser.

This $437\,\mathrm{bar}$ state corresponds to the reference setting adopted in the final simulation files, rather than an additional storage-folder identifier or a mandatory condition specified by SAE~J2601, it is used because the current project continues the work done for a manufacturer that required this specific state of the fluid.

So, the reference values in the final OpenFOAM dictionaries are $\rho_{\mathrm{H}_2}=33.77\,\mathrm{kg\,m^{-3}}$ and $\nu_{\mathrm{H}_2}=2.709\times10^{-7}\,\mathrm{m^2\,s^{-1}}$, with the corresponding secondary nitrogen settings documented in Chapter~\ref{ch:doe}.

At this state, Hydrogen is a dense, high-pressure gas and its density cannot reliably be inferred from the ideal-gas equation alone. Both $\rho(p,T)$ and the dynamic viscosity $\mu(p,T)$ should be determined using a consistent real-fluid database or correlation, such as NIST REFPROP, before assigning the constant kinematic viscosity $\nu=\mu/\rho$ and the reference density used by OpenFOAM's force calculation \cite{nist_refprop,nist_refprop_online}. It is important ot underline that the reference pressure is \emph{not} imposed by setting the kinematic outlet pressure to zero in the incompressible solver but is represented indirectly by selecting the fluid properties $\rho(p,T)$ $\And$ $\mu(p,T)$.

\section{Operating Principle of a Passive Excess-Flow Valve}
\label{sec:efv_principle}

An excess-flow valve is intended to shut off or restrict gas flow automatically when the flow exceeds a design threshold \cite{iso19880_3}. The working principle of this device can be represented as a balance between the hydrodynamic load acting on a movable element and the force exerted by a spring which together maintain the EFV in its normal, open position \cite{honeywell_efv}. The following explanation gives a physical model useful for interpreting the present CFD study; exact component details and reset mechanisms depend on the particular valve construction.


\begin{figure}[htbp]
\centering
\resizebox{\textwidth}{!}{% Conceptual schematic only; actual geometry and stroke depend on EFV design.
\begin{tikzpicture}[>=Latex,line width=0.75pt,
 lab/.style={font=\footnotesize,align=center},
 flow/.style={->,line width=1.1pt}]
\begin{scope}
\node[font=\small\bfseries] at (3,3.7) {(a) Normal-flow condition};
\draw[thick] (0,0.45)--(6.1,0.45) (0,2.55)--(6.1,2.55);
\draw[thick] (0,0.45)--(0,2.55) (6.1,0.45)--(6.1,2.55);
\draw[fill=black!22,thick] (2.6,1.02) rectangle (3.15,1.98);
\draw[fill=black!40] (4.52,1.05) rectangle (4.72,1.95);
\draw[thick] (3.15,1.5)--(3.4,1.5)--(3.58,1.78)--(3.82,1.22)--(4.05,1.78)--(4.3,1.22)--(4.52,1.5);
\draw[flow] (-0.45,2.22)--(1.7,2.22);
\draw[flow] (-0.45,0.78)--(1.7,0.78);
\draw[flow] (3.3,2.22)--(5.7,2.22);
\draw[flow] (3.3,0.78)--(5.7,0.78);
\node[lab] at (1.05,2.87) {Inlet flow};
\node[lab] at (2.78,-0.13) {Obturator};
\node[lab] at (4.7,-0.13) {Seat};
\node[lab] at (5.24,1.5) {Outlet};
\draw[->] (2.7,2.82)--(3.25,2.82) node[midway,above,font=\scriptsize] {$F_{\mathrm{fluid},x}$};
\draw[->] (4.3,2.82)--(3.75,2.82) node[midway,above,font=\scriptsize] {$F_{\mathrm{spring}}$};
\end{scope}
\begin{scope}[xshift=7.5cm]
\node[font=\small\bfseries] at (3,3.7) {(b) Excess-flow tendency};
\draw[thick] (0,0.45)--(6.1,0.45) (0,2.55)--(6.1,2.55);
\draw[thick] (0,0.45)--(0,2.55) (6.1,0.45)--(6.1,2.55);
\draw[fill=black!22,thick] (3.93,1.02) rectangle (4.48,1.98);
\draw[fill=black!40] (4.52,1.05) rectangle (4.72,1.95);
\draw[thick] (4.48,1.5)--(4.58,1.5);
\draw[flow] (-0.5,2.19)--(2.3,2.19);
\draw[flow] (-0.5,0.81)--(2.3,0.81);
\draw[flow] (4.82,2.19)--(5.7,2.19);
\draw[flow] (4.82,0.81)--(5.7,0.81);
\draw[->,thick] (3.10,1.5)--(3.78,1.5);
\node[lab] at (1.0,2.87) {Increased inlet flow};
\node[lab] at (3.95,-0.13) {Obturator near seat};
\node[lab] at (5.35,1.5) {Outlet};
\draw[->,thick] (3.15,2.82)--(4.05,2.82) node[midway,above,font=\scriptsize] {$F_{\mathrm{fluid},x}$};
\draw[->] (5.25,2.82)--(4.65,2.82) node[midway,above,font=\scriptsize] {$F_{\mathrm{spring}}$};
\end{scope}
\draw[->] (5.5,-0.65)--(7.4,-0.65) node[midway,below,font=\scriptsize] {$+x$: closure direction};
\end{tikzpicture}
}
\caption{Conceptual spring loaded EFV states: (a) normal flow, with the obturator away from the seat, and (b) increased hydrodynamic loading tending to move it towards closure \cite{iso19880_3,white2016}.}
\label{fig:efv_open_closing}
\end{figure}

\subsection{Normal-flow condition}

During normal operation, the movable element remains away from its shut-off seat, leaving flow passages through or around the obturator. The resulting total force is obtained from integration over the wetted obturator surface $S_{\mathrm{obt}}$:
\begin{equation}
\mathbf{F}_{\mathrm{fluid}}=
\int_{S_{\mathrm{obt}}}\left[-p\,\mathbf{n}+\boldsymbol{\tau}\cdot\mathbf{n}\right]\mathrm{d}S,
\label{eq:efv_surface_force}
\end{equation}
where $\mathbf{n}$ is the outward surface normal of the solid into the fluid, $p$ is static pressure and $\boldsymbol{\tau}$ is the viscous stress tensor; the pressure-force sign depends on the chosen normal orientation \cite{white2016,moukalled2016}. Since the obturator will move along the valve axis, the axial component of the applied forces will be the quantity most directly related to the tendency to close or open the valve.

A well designed EFV should exhibit an acceptable pressure loss in this state, since a restrictive geometry may produce greater fluid load, but also increase pumping or compression demand, interfering with the flow control margin of the refueling system. In the work executed, to account for these physical concepts both the pressure difference and obturator load rather are evaluated and compared.

\subsection{Excess-flow condition and motion towards closure}

Inside the valve an increased flow demand, rises the velocity in the open passages. The distribution of pressure and shear over the obturator cosequently changes, potentially increasing the hydrodynamic force directed towards the seat. To assess this phenomena, a simplified one-dimensional dynamic balance may be written as:
\begin{equation}
m\ddot{x}=
F_{\mathrm{fluid},x}(x,\dot{x},Q)
-F_0-kx-F_{\mathrm{friction}}-F_{\mathrm{other}},
\label{eq:efv_force_balance}
\end{equation}
where positive displacement $x$ is defined here as the direction of closure, $\dot{x}$ is the displacement time derivative, $F_0$ is the preload acting on the spring, $k$ is the spring stiffness, $F_{\mathrm{friction}}$ is the fricion force and the remaining term represent any additional relevant mechanical load. This expression is a conceptual balance, not a calibrated motion model for the present valve along the direction that will be used as the streamwise $x$ axis used to report CFD results.


\begin{figure}[htbp]
\centering
\resizebox{0.95\textwidth}{!}{% Conceptual free-body model and CFD outputs (same x-axis as the text).
\begin{tikzpicture}[>=Latex,line width=0.8pt,
  textnode/.style={font=\small,align=center}]
\draw[fill=black!18,thick] (3.5,0) rectangle (4.1,2.0);
\node[textnode] at (3.8,-0.43) {Movable obturator};
\draw[->,line width=1.2pt] (0.3,1.35)--(3.1,1.35)
 node[midway,above,textnode] {$F_{\mathrm{fluid},x}$};
\draw[->,line width=1.2pt] (7.0,1.35)--(4.5,1.35)
 node[midway,above,textnode] {$F_0+kx$};
\draw[densely dotted] (0.0,0.75)--(7.0,0.75);
\node[textnode] at (0.65,0.34) {$p_{\mathrm{in}}$};
\node[textnode] at (6.40,0.34) {$p_{\mathrm{out}}$};
\draw[->,thick] (3.7,-1.1)--(5.3,-1.1)
 node[midway,below,font=\footnotesize] {$+x$ towards seat};
\node[draw,rounded corners,inner sep=5pt,fill=black!3,textnode] at (3.6,-2.45)
 {$m\ddot{x}=F_{\mathrm{fluid},x}-F_0-kx-F_{\mathrm{other}}$};
\node[textnode] at (3.6,-3.2)
 {Fixed-obturator CFD provides a load at prescribed position,\\not the transient displacement $x(t)$.};
\end{tikzpicture}
}
\caption{Free-body diagram and governing equation for a spring loaded fixed obturator like the CFD geometry \cite{white2016,moukalled2016}.}
\label{fig:efv_force_balance}
\end{figure}

Closure becomes mechanically possible when the flow-induced load overcomes the effective opposing spring forces at the relevant position. Once the element approaches its seat, the effective passage area, local flow field and pressure distribution change. Consequently, the actual trip point is coupled to obturator displacement \cite{white2016,moukalled2016}.
A single CFD obturator position doesn't fully evaluates the behavior of the valve during closure but it is useful for the evaluation of the operational behavior at regime. 

\subsection{Shut-off, residual flow and resetting}

Depending on the valve design, seating may almost stop the gas flow or leave a calibrated residual leakage path. Some EFVs use a bleed-by passage to equalise pressure after the downstream fault has been removed, allowing the spring to return the element to the open position. Other designs require a deliberate manual reset \cite{honeywell_efv,linde_efv,teco_efv}. In a Refueling installation these details affect maintenance, fault recovery and the coordination between passive and actively controlled safety devices.

\section{Design parameters and adopted tradeoffs}
\label{sec:efv_design_parameters}

For an obturator with $n$ nominally circular holes of diameter $d$, the total geometric passage area is
\begin{equation}
A_{\mathrm{holes}}=n\frac{\pi d^2}{4}.
\label{eq:efv_open_area}
\end{equation}
At a prescribed flow rate, changing $n$ or $d$ alters the average velocity required through the passages but is important to underline that contraction and local separation can raise velocities above a simple estimation. Accordingly, the pressure field and the axial hydrodynamic force depend not only on the total area but also on hole distribution, surrounding wall shape and local flow structure \cite{crane410,white2016}. This implies that studying hole count and hole diameter should be separate.

A convenient nominal description of the pressure loss is
\begin{equation}
P=\frac{\Delta p}{\tfrac{1}{2}\rho U_{\mathrm{in}}^2}
\label{eq:efv_loss_coefficient}
\end{equation}
where $U_{\mathrm{in}}$ is the reference inlet velocity. Likewise, a dimensionless axial-force coefficient can be formed using a specified reference area $A_{\mathrm{ref}}$:
\begin{equation}
C_{F,x}=\frac{F_x}{\tfrac{1}{2}\rho U_{\mathrm{in}}^2 A_{\mathrm{ref}}}
\label{eq:efv_force_coefficient}
\end{equation}
These coefficients help distinguish geometric effects from changes in operating flow, allowing so to compare the different simulations between each other \cite{white2016}.

\begin{table}[htbp]
\centering
\small
\begin{tabular}{@{}p{0.25\textwidth}p{0.35\textwidth}p{0.35\textwidth}@{}}
\toprule
\textbf{Design point} & \textbf{Fluid dynamic implication} & \textbf{Engineering implication}\\
\midrule
Holes number $\And$ size & Changes local restriction, velocity and force  & Balances normal-flow loss with available actuation load\\
Spring $K$ $\And$ $F_0$ & Determines opposing mechanical load & Sets response threshold and reset capability\\
Upstream gas state & Sets density, viscosity and compression possibility & Changes pressure drop and force scaling\\
Obturator $x$ position & Alters effective passage area and pressure distribution & Couples steady forces to the closing phase\\
\bottomrule
\end{tabular}
\caption{Principal design variables affecting a spring loaded EFV}
\label{tab:efv_design_tradeoffs}
\end{table}

\section{CFD Study within the Refueling Application}
\label{sec:efv_cfd_scope}

The present numerical investigation isolates a fixed obturator in a three-dimensional internal-flow domain. For each chosen gas, inlet velocity and geometry, it estimates the pressure loss, the axial force component and related flow field quantities. 

These outputs are directly relevant to the fluid dynamic part of an EFV design assessment and can inform selection among candidate hole configurations.\\

The work does \emph{not} simulate a complete Refueling-station filling sequence. In particular, SAE~J2601 fuel-delivery temperatures and nominal pressure classes define an \emph{application context}, since there's no proof that a steady, constant density calculation reproduces every refueling condition, even if the fluid initial properties are the same \cite{sae_j2601}. To complete the design choices on a EFV inside these cùkind of system  would also require to address the pressure and temperature rating, sealing behaviour, materials compatibility, transient response and the applicable test procedures of ISO~19880-3 \cite{iso19880_3}. \\

The following results of this study so, will be only one of the many elements needed to make the system designer able to choose properly the valve to be installed.
